%==============================================================
%  lecons/analyse/lebesgue.tex — Intégrale de Lebesgue
%==============================================================

\lecontitle{Intégrale de Lebesgue}{Analyse réelle — Mesure et intégration}

%=============================================================
%  § 1 — CONSTRUCTION
%=============================================================
\secnum{navyblue}{1}{De la mesure à l'intégrale}

\begin{tcolorbox}[thmbox]
  \textbf{Idée directrice.}\enspace%
  \index{intégrale de Lebesgue}
  Là où Riemann découpe l'\emph{abscisse} (subdivisions de $[a,b]$),
  Lebesgue découpe l'\emph{ordonnée} : il regroupe les points $t$ selon
  la valeur $f(t)$, puis pèse chaque tranche par une \emph{mesure}.
  La mesure de Lebesgue $\lambda$ généralise la longueur~: $\lambda([a,b])=b-a$.

  \medskip
  \textbf{Étape 1 — Tribu et mesure.}\enspace%
  \index{tribu borélienne}\index{mesure de Lebesgue}
  On dispose d'une tribu $\mathcal{B}(\mathbb{R})$ (borélienne) et d'une mesure
  $\lambda : \mathcal{B}(\mathbb{R}) \to [0,+\infty]$, $\sigma$-additive,
  invariante par translation, normalisée par $\lambda([0,1])=1$.

  \medskip
  \textbf{Étape 2 — Fonctions étagées.}\enspace%
  Une fonction $\varphi = \sum_{k=1}^N \alpha_k \mathbf{1}_{A_k}$ avec
  $\alpha_k \geq 0$ et $A_k$ mesurables disjoints est dite
  \emph{étagée positive}. On pose
  \[
    \int \varphi\,\mathrm{d}\lambda \;=\; \sum_{k=1}^N \alpha_k\,\lambda(A_k)
    \;\in\; [0,+\infty].
  \]

  \medskip
  \textbf{Étape 3 — Fonctions mesurables positives.}\enspace%
  Toute fonction mesurable $f \geq 0$ est limite croissante d'une suite
  $(\varphi_n)$ de fonctions étagées positives. On définit
  \[
    \int f\,\mathrm{d}\lambda \;=\; \sup\!\left\{\int \varphi\,\mathrm{d}\lambda
    \;\Big|\; \varphi\ \text{étagée},\ 0\leq\varphi\leq f\right\} \;\in\; [0,+\infty].
  \]

  \medskip
  \textbf{Étape 4 — Fonctions de signe quelconque.}\enspace%
  Pour $f$ mesurable, on décompose $f = f^+ - f^-$ avec $f^{\pm} \geq 0$.
  Si $\int f^+ + \int f^- < +\infty$, $f$ est dite \emph{Lebesgue-intégrable}
  (notation $f \in L^1$) et l'on pose
  \[
    \int f\,\mathrm{d}\lambda \;=\; \int f^+\,\mathrm{d}\lambda
    \;-\; \int f^-\,\mathrm{d}\lambda.
  \]
\end{tcolorbox}

\begin{significationintuitive}
Pour calculer l'aire sous $f$, Lebesgue fait des tranches horizontales :
pour chaque hauteur $y$, on mesure l'ensemble $\{t : f(t) > y\}$. On
intègre ensuite ces mesures :
$\int f = \int_0^{+\infty} \lambda(\{f > y\})\,\mathrm{d}y$ pour $f \geq 0$.
Cette construction permet d'intégrer toute fonction mesurable positive,
y compris très irrégulière, et donne un théorème de convergence simple :
si $f_n \to f$ presque partout avec une domination intégrable, alors
$\int f_n \to \int f$. Riemann ne dispose pas d'un tel théorème.
\end{significationintuitive}

\decsep{}

%=============================================================
%  § 2 — THÉORÈMES FONDAMENTAUX
%=============================================================
\secnum{navyblue}{2}{Théorèmes de convergence}

\begin{tcolorbox}[thmbox]
  \textbf{Théorème 1 (convergence monotone, Beppo Levi).}\enspace%
  \index{convergence monotone (théorème de)}
  Soit $(f_n)$ une suite croissante de fonctions mesurables positives,
  convergeant simplement vers $f$. Alors
  \[
    \int f_n\,\mathrm{d}\lambda \;\xrightarrow[n\to+\infty]{}\;
    \int f\,\mathrm{d}\lambda \quad \text{(dans }[0,+\infty]\text{)}.
  \]

  \medskip
  \textbf{Théorème 2 (lemme de Fatou).}\enspace\index{Fatou (lemme de)}
  Pour toute suite $(f_n)$ de fonctions mesurables positives,
  \[
    \int \liminf_n f_n\,\mathrm{d}\lambda
    \;\leq\; \liminf_n \int f_n\,\mathrm{d}\lambda.
  \]

  \medskip
  \textbf{Théorème 3 (convergence dominée, Lebesgue).}\enspace%
  \index{convergence dominée (théorème de)}
  Soit $(f_n)$ mesurables convergeant simplement (ou p.p.) vers $f$.
  On suppose qu'il existe $g \in L^1$ positive telle que $|f_n| \leq g$
  pour tout $n$ (\emph{hypothèse de domination}). Alors $f \in L^1$ et
  \[
    \int f_n\,\mathrm{d}\lambda \;\xrightarrow[n\to+\infty]{}\;
    \int f\,\mathrm{d}\lambda,
    \qquad
    \int |f_n - f|\,\mathrm{d}\lambda \to 0.
  \]

  \medskip
  \textbf{Théorème 4 (Fubini--Tonelli).}\enspace%
  \index{Fubini--Tonelli (théorème de)}
  Pour $f : \mathbb{R}^2 \to \mathbb{R}$ mesurable~:
  \begin{itemize}
    \item Si $f \geq 0$ (Tonelli) : les trois intégrales
    \[
      \iint f\,\mathrm{d}\lambda^2,\quad
      \int\!\!\!\int f\,\mathrm{d}x\,\mathrm{d}y,\quad
      \int\!\!\!\int f\,\mathrm{d}y\,\mathrm{d}x
    \]
    sont toutes égales dans $[0,+\infty]$.
    \item Si $\iint |f| < +\infty$ (Fubini) : $f \in L^1(\mathbb{R}^2)$ et
    les intégrales itérées dans n'importe quel ordre coïncident.
  \end{itemize}
\end{tcolorbox}

\rubriq{Lien avec Riemann}

\begin{tcolorbox}[thmbox]
  \textbf{Théorème (Riemann $\subset$ Lebesgue).}\enspace%
  Soit $f:[a,b]\to\mathbb{R}$ bornée. Alors
  \begin{itemize}
    \item $f$ est Riemann-intégrable si et seulement si l'ensemble de ses
          points de discontinuité est de mesure nulle (\emph{critère de
          Lebesgue}\index{critère de Lebesgue})~;
    \item dans ce cas, $f$ est Lebesgue-intégrable et les deux intégrales
          coïncident :
          $\displaystyle\int_a^b f(t)\,\mathrm{d}t = \int_{[a,b]} f\,\mathrm{d}\lambda.$
  \end{itemize}
  \textbf{Cas des intégrales impropres.}\enspace Pour $f$ localement
  Riemann-intégrable sur un intervalle, $f$ est Lebesgue-intégrable si et
  seulement si l'intégrale impropre de $|f|$ converge, et les deux
  intégrales coïncident alors.
  Les intégrales semi-convergentes (type $\int_0^{+\infty}\sin(t)/t\,\mathrm{d}t$)
  sont définies par Riemann mais pas par Lebesgue.
\end{tcolorbox}

\decsep{}

%=============================================================
%  § 3 — DÉMONSTRATIONS (esquissées)
%=============================================================
\secnum{navyblue}{3}{Éléments de démonstration}

\rubriq{Convergence monotone}

\begin{mdframed}[style=proofbar]
  Soit $f_n \uparrow f$ avec $f_n \geq 0$ mesurables. La suite
  $\bigl(\int f_n\bigr)$ est croissante majorée par $\int f$, donc admet
  une limite $L \leq \int f$.

  \medskip
  \textit{Inégalité inverse.}\enspace%
  Pour montrer $L \geq \int f$, on prend $\varphi$ étagée avec
  $0 \leq \varphi \leq f$, $0 < c < 1$, et $A_n = \{f_n \geq c\varphi\}$.
  Comme $f_n \uparrow f \geq \varphi > c\varphi$ là où $\varphi>0$,
  $A_n \uparrow$ l'ensemble entier. Par $\sigma$-additivité,
  $\int c\varphi\,\mathbf{1}_{A_n} \to \int c\varphi$. Mais
  $\int f_n \geq \int c\varphi\,\mathbf{1}_{A_n}$, donc à la limite
  $L \geq c\int\varphi$. Faisant $c \to 1^-$ puis $\sup_\varphi$ :
  $L \geq \int f$.\hfill$\blacksquare$
\end{mdframed}

\rubriq{Lemme de Fatou}

\begin{mdframed}[style=proofbar]
  Posons $g_n = \inf_{k\geq n} f_k$. La suite $(g_n)$ est croissante et
  $g_n \to \liminf f_n$. Par convergence monotone,
  $\int g_n \to \int \liminf f_n$. Comme $g_n \leq f_n$,
  $\int g_n \leq \int f_n$, donc à la limite inférieure :
  $\int \liminf f_n \leq \liminf \int f_n$.\hfill$\blacksquare$
\end{mdframed}

\rubriq{Convergence dominée}

\begin{mdframed}[style=proofbar]
  Soit $|f_n|\leq g \in L^1$ et $f_n \to f$ p.p. Alors $|f|\leq g$ p.p., donc
  $f \in L^1$. Les fonctions $g + f_n \geq 0$ et $g - f_n \geq 0$ sont
  positives ; on leur applique le lemme de Fatou.

  \smallskip
  \textit{Premier terme.}\enspace
  Comme $g + f_n \to g + f$, Fatou donne
  \[
    \int g + \int f
    \;=\; \int (g + f)
    \;\leq\; \liminf_n \int (g + f_n)
    \;=\; \int g + \liminf_n \int f_n,
  \]
  d'où, $\int g$ étant fini, $\displaystyle\int f \leq \liminf_n \int f_n$.

  \smallskip
  \textit{Second terme.}\enspace
  De même, $g - f_n \to g - f$ et Fatou donne
  $\int g - \int f \leq \liminf_n \int (g - f_n)$. Or
  $\liminf_n\!\big({-}\!\int f_n\big) = -\limsup_n \int f_n$, donc
  $\liminf_n \int (g - f_n) = \int g - \limsup_n \int f_n$. En simplifiant
  par $\int g$ on obtient
  $\displaystyle \limsup_n \int f_n \leq \int f$.

  \smallskip
  Les deux inégalités donnent
  $\displaystyle \limsup_n \int f_n \leq \int f \leq \liminf_n \int f_n$,
  d'où $\int f_n \to \int f$.

  \medskip
  Pour la convergence dans $L^1$ : appliquer Fatou à $|f_n - f| \leq 2g$
  avec $\liminf |f_n - f| = 0$.\hfill$\blacksquare$
\end{mdframed}

\rubriq{Comparaison Riemann/Lebesgue (sens facile)}

\begin{mdframed}[style=proofbar]
  Soit $f:[a,b]\to\mathbb{R}$ continue. Pour une subdivision régulière
  $\sigma_n$ de pas $1/n$, les sommes de Darboux
  $s(f,\sigma_n) \to \int_{[a,b]} f\,\mathrm{d}\lambda$ par convergence
  dominée appliquée aux fonctions étagées
  $\varphi_n = \sum_i m_i^{(n)} \mathbf{1}_{[x_{i-1},x_i[}$ : elles convergent
  simplement vers $f$ (continuité de $f$) tout en restant majorées par la
  constante $\sup_{[a,b]}|f|$, intégrable sur $[a,b]$. (On évite ici la
  convergence monotone, car les $\sigma_n$ ne sont pas emboîtées et la suite
  $(\varphi_n)$ n'est pas croissante.) De même
  $S(f,\sigma_n) \to \int_{[a,b]} f\,\mathrm{d}\lambda$. Comme $f$ est
  Riemann-intégrable, $s, S \to \int_a^b f$ (Riemann). Les deux limites
  coïncident.\hfill$\blacksquare$
\end{mdframed}

\decsep{}

%=============================================================
%  § 4 — EXEMPLES, CONTRE-EXEMPLES, EXERCICES
%=============================================================
\secnum{exgreen}{4}{Exemples, contre-exemples et exercices}

\rubriq{Exemples}

\begin{mdframed}[style=exbar]
  \textbf{1. Dirichlet est Lebesgue-intégrable.}\enspace%
  $\mathbf{1}_{\mathbb{Q}\cap[0,1]}$ vaut $0$ sauf sur un ensemble dénombrable,
  donc de mesure nulle :
  $\int_{[0,1]} \mathbf{1}_{\mathbb{Q}}\,\mathrm{d}\lambda = 0$,
  alors que Riemann ne sait pas l'intégrer.

  \smallskip\noindent
  \textbf{2. Limite simple non dominée.}\enspace%
  $f_n(t) = n\,\mathbf{1}_{[0,1/n]}(t)$ vérifie $f_n \to 0$ p.p.\ mais
  $\int f_n = 1$ pour tout $n$. La domination $|f_n| \leq g$ est violée
  (la dominante minimale $\sup_n f_n\notin L^1$).

  \smallskip\noindent
  \textbf{3. Échange dérivation-intégrale.}\enspace%
  Pour $F(x) = \int_\mathbb{R} f(x,t)\,\mathrm{d}\lambda(t)$ avec $f$ dérivable
  en $x$ et $|\partial_x f(x,t)| \leq g(t) \in L^1$ uniformément en $x$,
  la convergence dominée donne $F'(x) = \int \partial_x f(x,t)\,\mathrm{d}\lambda(t)$.

  \smallskip\noindent
  \textbf{4. Calcul de $\int_0^{+\infty}\frac{\sin t}{t}\,\mathrm{d}t$.}\enspace%
  Cette intégrale vaut $\pi/2$ au sens de Riemann (limite de $\int_0^A$),
  mais $\int_0^{+\infty}|\sin t/t| = +\infty$, donc \emph{pas}
  Lebesgue-intégrable.

  \smallskip\noindent
  \textbf{5. Fubini.}\enspace%
  $\displaystyle\int_0^1\!\!\int_0^1 \frac{x^2 - y^2}{{(x^2+y^2)}^2}\,\mathrm{d}y\,\mathrm{d}x
  = \frac{\pi}{4}$ alors que l'intégrale dans l'autre ordre vaut $-\pi/4$ :
  l'hypothèse $\iint |f| < +\infty$ est violée, Fubini ne s'applique pas.
\end{mdframed}

\vspace{6pt}
\rubriq{Contre-exemples et subtilités}

\begin{mdframed}[style=cexbar]
  \textbf{Mesure nulle $\neq$ ensemble petit.}\enspace%
  L'ensemble de Cantor triadique est non dénombrable, mais de mesure nulle :
  $\lambda(\mathcal{C}) = 0$. La fonction indicatrice $\mathbf{1}_\mathcal{C}$
  est d'intégrale nulle.

  \smallskip\noindent
  \textbf{Convergence p.p.\ sans convergence dans $L^1$.}\enspace%
  $f_n = n\,\mathbf{1}_{[0,1/n]} \to 0$ p.p., mais $\|f_n\|_1 = 1 \not\to 0$.
  Sans hypothèse de domination, on ne contrôle pas l'intégrale de la limite.

  \smallskip\noindent
  \textbf{Lebesgue ne « voit » pas la semi-convergence.}\enspace%
  $\int_0^{+\infty}\sin(t)/t\,\mathrm{d}t$ existe au sens de Riemann impropre
  ($=\pi/2$) mais pas au sens de Lebesgue. Lebesgue est strictement plus
  exigeant pour les intégrales infinies (il demande absolue convergence).

  \smallskip\noindent
  \textbf{Une fonction Lebesgue-intégrable peut être nulle part Riemann-intégrable.}\enspace%
  $\mathbf{1}_{\mathbb{Q}\cap[0,1]}$ est partout discontinue : aucun
  sous-intervalle ne porte $f$ Riemann-intégrable.
\end{mdframed}

\vspace{6pt}
\exolabel

\begin{mdframed}[style=exobar]
  \begin{enumerate}
    \item Soit $f_n(t) = \dfrac{n}{1+n^2 t^2}$ sur $[0,1]$. Montrer que
          $f_n \to 0$ p.p.\ mais $\int_0^1 f_n \not\to 0$. Identifier la
          masse de Dirac qui se cache derrière.

    \item Montrer que $\displaystyle\int_0^{+\infty}\frac{\mathrm{d}t}{1+t^n}
          \xrightarrow[n\to+\infty]{} 1$ par convergence dominée.

    \item Démontrer la formule
          $\displaystyle\int_0^{+\infty} e^{-t}t^{s-1}\,\mathrm{d}t = \Gamma(s)$,
          de classe $\mathcal{C}^\infty$ en $s>0$, par dérivation
          sous le signe intégral.

    \item Soit $f \in L^1(\mathbb{R})$. Montrer la \emph{continuité dans $L^1$}~:
          $\|f(\cdot - h) - f\|_1 \to 0$ quand $h\to 0$.
          (\textit{Indication}~: densité des fonctions $\mathcal{C}_c$ dans $L^1$.)

    \item Établir la formule
          $\displaystyle\int_0^{+\infty}\frac{\ln t}{1+t^2}\,\mathrm{d}t = 0$
          par changement $t \mapsto 1/t$, ou via Fubini en écrivant
          $\ln t = \displaystyle\int_0^{+\infty}\Bigl(\frac{1}{1+u}-\frac{1}{t+u}\Bigr)\mathrm{d}u$.

    \item Soit $A \subset [0,1]$ borélien avec $\lambda(A) = 1$. Montrer
          que $A$ est dense dans $[0,1]$.

    \item (\textit{Ensemble de Cantor « gras »}) Construire un sous-ensemble
          fermé $K \subset [0,1]$ d'intérieur vide et de mesure $\lambda(K) = 1/2$.
          Montrer que l'ensemble des points de discontinuité de $\mathbf{1}_K$
          est exactement $K$, puis en déduire, via le critère de Lebesgue, que
          $\mathbf{1}_K$ n'est \emph{pas} Riemann-intégrable bien que $K$ soit
          d'intérieur vide.

    \item Soit $f_n \to f$ dans $L^1$. Montrer qu'il existe une sous-suite
          $(f_{n_k})$ convergeant vers $f$ presque partout.
          (\textit{Indication}~: choisir $n_k$ tel que $\|f_{n_k}-f\|_1 \leq 2^{-k}$
          et appliquer Borel--Cantelli.)
  \end{enumerate}
\end{mdframed}

\leconfooter{H.\ Lebesgue, \textit{Intégrale, longueur, aire} (1902) — É.\ Borel — G.\ Vitali}
